\clearpage % Rozdziały zaczynamy od nowej strony.
\section{Podsumowanie i dalsze kierunki badań}

Niniejsza praca zbadała, w jakim stopniu agenty oparte na dużych modelach językowych mogą autonomicznie generować konfiguracje Infrastructure as Code — pliki Dockerfile i manifesty Kubernetes — dla rzeczywistych aplikacji webowych. Eksperymenty objęły pięć pytań badawczych dotyczących skuteczności funkcjonalnej, wpływu złożoności, jakości artefaktów, powtarzalności procesu oraz odporności na manipulację kontekstem.

\subsection{Odpowiedzi na pytania badawcze}

\textbf{RQ1.} Agenty LLM potrafią autonomicznie wygenerować działające konfiguracje IaC dla większości prostych i umiarkowanie złożonych aplikacji webowych. Granica ich możliwości leży nie w znajomości składni, lecz w zdolności do przewidywania zachowania środowiska runtime bez możliwości jego obserwacji.

\textbf{RQ2.} Złożoność architektury wielousługowej stanowi barierę, której obecne agenty nie przekraczają niezawodnie. Wraz ze wzrostem liczby współzależnych serwisów skuteczność wyraźnie spada — agent generuje pliki lokalnie poprawne, które nie tworzą spójnej całości. Kierunek tego trendu jest wiarygodny, jednak ze względu na małą liczbę kontrolowanych repozytoriów PoC (n=5) nie można precyzyjnie określić kształtu krzywej ani wskazać ilościowego „punktu załamania".

\textbf{RQ3.} Domyślna jakość generowanych artefaktów jest niewystarczająca do zastosowań produkcyjnych. Rozszerzenie instrukcji systemowej o wytyczne jakościowe drastycznie poprawia wyniki, co wskazuje, że modele posiadają potrzebną wiedzę, lecz wymagają jawnego nakierowania.

\textbf{RQ4.} Proces agentowy nie jest w pełni powtarzalny nawet przy deterministycznych ustawieniach modelu. Zmienność wynika z wieloetapowej architektury agenta i narasta wraz ze złożonością zadania.

\textbf{RQ5.} Agenty są podatne na manipulację dokumentacją repozytorium. Fałszywe wymagania osadzone w plikach tekstowych skutecznie wpływają na generowane konfiguracje bez konieczności exploitowania modelu, przy czym podatność jest silnie zróżnicowana między modelami.

\subsection{Odpowiedź na główne pytanie badawcze}

Niniejsza praca stawiała pytanie, czy agenty LLM są w stanie — bez ingerencji człowieka — wygenerować funkcjonalne, bezpieczne i zgodne z dobrymi praktykami konfiguracje wdrożeniowe dla rzeczywistych aplikacji webowych. Odpowiedź jest warunkowa: \textbf{tak, dla prostych aplikacji jednokomponentowych z czytelną strukturą kodu; nie, dla złożonych systemów wielousługowych oraz środowisk o podwyższonych wymaganiach bezpieczeństwa lub z niezaufaną dokumentacją}.

Przy prostych aplikacjach agenty osiągają ponad 60\% end-to-end success rate i mogą generować konfiguracje nadające się do wdrożenia bez interwencji człowieka. Przy rosnącej złożoności skuteczność spada drastycznie (do 0\% dla najbardziej złożonych przypadków), domyślna jakość artefaktów jest niewystarczająca bez odpowiedniego prompt engineeringu, a podatność na manipulację kontekstem czyni w pełni autonomiczne zastosowanie ryzykownym wszędzie tam, gdzie repozytorium może zawierać niezaufaną dokumentację.

Warto jednak zastrzec, że obserwowane ograniczenia nie są koniecznie ograniczeniami samych modeli językowych — lecz architektury systemu agentowego, w którym są osadzone. Kluczowym czynnikiem jest brak pętli sprzężenia zwrotnego z runtime: agent generuje konfiguracje statycznie, nie mogąc obserwować ani reagować na błędy budowania kontenera czy awarie wdrożeniowe. Iteracyjna korekta oparta na faktycznym wyjściu systemu — jak w projektach takich jak Repo2Run — mogłaby istotnie podnieść skuteczność przy niezmienionych modelach bazowych. W tym sensie wyniki pracy charakteryzują konkretną architekturę agentową, a nie granicę możliwości LLM jako takich.

\subsection{Kluczowe wnioski}

Wyniki wskazują, że agenty LLM stanowią realne, lecz ograniczone narzędzie do automatyzacji IaC. Są użyteczne jako asystent inżyniera przy prostych, jednokomponentowych aplikacjach — skracając czas tworzenia szkieletu konfiguracji — ale nie zastępują człowieka w scenariuszach złożonych lub wymagających wysokiego poziomu bezpieczeństwa. Kluczowym odkryciem jest asymetria między tym, co model \textit{potrafi} (zna składnię, zna dobre praktyki), a tym, co \textit{robi bez nakierowania} (generuje konfiguracje poniżej standardu, ulega manipulacji). Ta asymetria oznacza, że jakość systemu agentowego zależy w dużej mierze od jakości jego projektu — instrukcji, zestawu narzędzi i mechanizmów walidacji — a nie tylko od wyboru modelu bazowego.

Żaden z badanych modeli nie dominuje jednocześnie we wszystkich wymiarach: skuteczności funkcjonalnej, powtarzalności i odporności na manipulację. Wybór modelu powinien być podyktowany profilem ryzyka konkretnego wdrożenia.

\subsection{Implikacje praktyczne}

Bezpieczne zastosowanie agentów LLM do generowania IaC wymaga spełnienia kilku warunków jednocześnie:

\begin{itemize}
    \item \textbf{Ograniczony kontekst wejściowy} — agent powinien mieć dostęp wyłącznie do plików kodu źródłowego, bez dokumentacji i historii repozytorium, co istotnie ogranicza podatność na manipulację (RQ5).
    \item \textbf{Rozbudowana instrukcja systemowa} — prompt zawierający konkretną, sprawdzalną checklistę dobrych praktyk redukuje liczbę ostrzeżeń statycznych o ponad 79\% i trzykrotnie zwiększa odsetek konfiguracji wolnych od ostrzeżeń (RQ3).
    \item \textbf{Obowiązkowa walidacja statyczna} — narzędzia takie jak Hadolint i Kube-linter powinny być elementem pipeline CI/CD dla każdej konfiguracji IaC generowanej przez LLM.
    \item \textbf{Ludzki przegląd konfiguracji bezpieczeństwa} — artefakty zawierające security contexts, uprawnienia i sekrety wymagają weryfikacji człowieka przed wdrożeniem, niezależnie od zastosowanego modelu.
    \item \textbf{Ograniczenie do prostych zastosowań} — dla aplikacji wielousługowych (trzy lub więcej wzajemnie zależnych komponentów) agenty nie nadają się do w pełni autonomicznej pracy; wartościowe jest wówczas generowanie poszczególnych komponentów z oddzielną weryfikacją spójności całości.
\end{itemize}

\subsection{Ograniczenia pracy}
\label{sec:ograniczenia}

Wyniki dotyczą konkretnej implementacji systemu agentowego, określonego zestawu modeli komercyjnych i zbioru prostych aplikacji webowych w językach Python i Node.js. Nie można ich bezpośrednio uogólniać na inne technologie, modele open-source ani środowiska produkcyjne. Dodatkowym ograniczeniem jest brak pętli sprzężenia zwrotnego z rzeczywistym środowiskiem wykonawczym — architektura generująca konfiguracje statycznie może zaniżać wyniki względem systemów z iteracyjną korektą. Retrospektywna ocena pozostałych zagrożeń trafności (zależność od promptu i narzędzi agenta, brak walidacji runtime w RQ3 i RQ5, wpływ środowiska MicroK8s) znajduje się w sekcji~\ref{sec:retrospective} rozdziału~6.

Osobnym ograniczeniem dotyczy wniosków z RQ2. Pięć kontrolowanych repozytoriów PoC stanowi zbyt małą próbę, aby precyzyjnie opisać kształt zależności między złożonością a skutecznością — obserwowany monotonicznie malejący trend jest wiarygodny jako kierunek, lecz dokładna krzywa, a w szczególności lokalizacja ewentualnego punktu krytycznego, pozostaje nieokreślona. Dodatkowo, autor zaprojektował zarówno system testowy, jak i same repozytoria PoC, co stwarza ryzyko nieświadomego dobrania zadań do możliwości systemu; wyniki RQ2 należy zatem traktować jako wstępną ilustrację zjawiska, a nie jego pełną charakterystykę ilościową.

\subsection{Dalsze kierunki badań}

Wyniki pracy otwierają kilka kierunków dalszych badań:

\begin{itemize}
    \item \textbf{Architektura wieloagentowa z podziałem ról.} Zbadanie, czy separacja ról planowania, generowania i walidacji między wyspecjalizowanymi agentami redukuje zmienność i poprawia koordynację wielousługową.

    \item \textbf{Pętla sprzężenia zwrotnego z runtime.} Rozszerzenie systemu o iteracyjną korektę na podstawie wyników faktycznego uruchomienia kontenera, co może istotnie zwiększyć skuteczność etapu runtime.

    \item \textbf{Mechanizmy obronne przed manipulacją kontekstem.} Opracowanie i ocena technik wykrywania lub neutralizowania repo poisoning — weryfikacja spójności między kodem a dokumentacją, analiza anomalii lub ograniczenie uprawnień informacyjnych agenta.

    \item \textbf{Fine-tuning i szerszy zestaw modeli.} Porównanie modeli ogólnych z modelami dostrojonymi na zbiorach konfiguracji IaC oraz weryfikacja wyników na modelach open-source i innych stosach technologicznych.

    \item \textbf{Integracja z praktykami CI/CD.} Ocena agentów LLM jako elementu automatycznego pipeline wdrożeniowego w rzeczywistych projektach, z pomiarem wpływu na czas i jakość konfiguracji.
\end{itemize}

Agenty LLM już dziś mogą pełnić rolę użytecznego asystenta inżyniera DevOps — pod warunkiem świadomego projektowania systemu, właściwej kontroli kontekstu i obowiązkowej walidacji artefaktów. Pełna autonomia, rozumiana jako zastąpienie człowieka w całym procesie tworzenia i wdrażania konfiguracji infrastrukturalnej, pozostaje celem wymagającym rozwiązania otwartych problemów koordynacji wielousługowej, powtarzalności procesu i odporności na manipulację.
